% Part: computability
% Chapter: computability-theory
% Section: no-universal-function

\documentclass[../../../include/open-logic-section]{subfiles}

\begin{document}

\olfileid{cmp}{thy}{nou}
\olsection{કોઈ સર્વવ્યાપી સંપૂર્ણ સંગણનીય વિધેય નથી}

આંશિક સંગણનીય વિધેયો માટે સર્વવ્યાપી આંશિક સંગણનીય વિધેય તો છે,
પરંતુ સંપૂર્ણ સંગણનીય વિધેયો માટે સર્વવ્યાપી હોય એવું સંપૂર્ણ
સંગણનીય વિધેય નથી.

\begin{thm}
\ollabel{thm:no-univ}
કોઈ સર્વવ્યાપી સંગણનીય વિધેય નથી. બીજા શબ્દોમાં, જો
$\fn{Un}'(k, x)$ એવું વિધેય હોય કે દરેક સંપૂર્ણ સંગણનીય વિધેય
$f(x)$ માટે એવી પ્રાકૃતિક સંખ્યા~$k$ હોય કે દરેક~$x$ માટે
$f(x) = \fn{Un}'(k,x)$, તો તે સંગણનીય નથી.
\end{thm}

\begin{proof}
આ સાદી વિકર્ણીકરણ સાબિતી છે: $\fn{Un}'(k,x)$ સંપૂર્ણ અને સંગણનીય
હોય, તો
\[
d(x) = \fn{Un}'(x, x) + 1
\]
પણ સંપૂર્ણ અને સંગણનીય થાય. પરંતુ વ્યાખ્યા મુજબ $d(k)$ અને
$\fn{Un}'(k,k)$ સમાન નથી. તેથી દરેક $k$ માટે $d(x)$ અને
$\fn{Un}'(k, x)$નાં મૂલ્યો ઓછામાં ઓછા એક~$x$ પર, એટલે $x = k$
પર, જુદાં પડે છે.
\end{proof}

\begin{explain}
ઉપરનું \olref[uni]{thm:univ-comp} બતાવે છે કે સર્વવ્યાપી વિધેયને
આંશિક રહેવા દઈએ તો આ વિકર્ણીકરણ દલીલથી બચી શકાય છે. એટલે
$\fn{Un}$ સંપૂર્ણ સંગણનીય વિધેયો માટે પણ સર્વવ્યાપી છે, પરંતુ તે
પોતે સંપૂર્ણ નથી. આંશિક સ્થિતિમાં વિકર્ણીકરણ દલીલ ચાલતી નથી.
\end{explain}

\begin{prob}
  \olref{thm:no-univ}ની સાબિતીની વિકર્ણીકરણ દલીલ આંશિક
  સ્થિતિમાં કેમ ચાલતી નથી તે સમજવા, વિધેય
  $f(x) \simeq \fn{Un}(x,x)+1$ લો. શું તે આંશિક સંગણનીય છે?
  જો હા, તો તેનો કોઈ સૂચક~$e$ હશે, એટલે $f(x) \simeq
  \fn{Un}(e,x)$. તો~$f(e)$ વિશે શું કહી શકાય?
\end{prob}


\end{document}
